\section{논의}
\label{sec:discussion}

\begin{revision}

\subsection{검사 가능성과 도구의 역할}
조건별 의무와 해제 증거를 연결하면, 새로운 근거 B가 앞선 A를 해제하는지와 A가 미해제인지 미상인지를 구분할 수 있다. 최초 문제 사건ㆍ관련 의무ㆍ사유를 원문 구절에 연결하는 출력은 수정할 주석이나 추가 증거가 필요한 항목을 검토하는 데 쓸 수 있다. 그러나 작성 프레임의 정확성을 모델 검사가 보장하지는 않는다. 잘못 연결한 조건 ID나 누락된 의무는 별도 의미 검토가 필요하다.

이전 이력 창과 전체 이력의 비교는 기억 길이가 일부 판정에 영향을 주었음을 보여 준다. FSM의 144/144와 보완 모델의 동일 사례 144/144는 이 범위의 논리를 단순 상태 기계로도 실행할 수 있다는 근거다. UPPAAL의 역할은 표현한 상태ㆍ조건 속성을 명시적으로 질의하고 진단 경로와 모델 파일을 보존하는 데 있다. 정확도 우위나 UPPAAL만의 필요성을 입증하지 않는다.

\subsection{주석 검토에서 구분할 조치}
확정 모순은 미해제가 확인된 의무와 행동을 함께 검토한다. 근거 부족은 A의 해제 여부를 B로 대신 채우지 않고 추가 주석/관측을 요청할 대상으로 둔다. 양립 판정도 명시된 의무에 대해서만 유효하다. P4는 주석을 수정할 근거가 아니라 검사 상태 처리의 고장을 점검할 근거다. 따라서 모든 경고를 자연 오류라고 하거나 자료를 자동 삭제하는 절차로 해석하지 않는다.

\subsection{검토 기록의 활용과 보존}
검토자는 먼저 문제 사건과 관련 의무를 원문에서 함께 읽는다. A=F의 모순이면 A의 미해제 구절과 GO 구절이 같은 조건ㆍ대상에 대응하는지 확인하고, 잘못 연결했다면 프레임 연결을 수정할 대상으로 둔다. A가 미상인 근거 부족이면 존재하지 않는 해제 사실을 정상으로 채우지 않고 필요한 객체 관측이나 주석을 구분한다. A=T의 양립이면 해당 의무를 왜 해제했는지와 출발 근거 B를 각각 보존한다. 이처럼 같은 ``검토 필요'' 표시도 문장 수정과 증거 보충이라는 다른 조치로 이어진다.

이 과정에서 유지할 기록은 원문 구절ㆍ조건 IDㆍ사건 순서ㆍ모델/규칙 버전ㆍ질의ㆍ최초 위치ㆍ사유다. 규칙이나 속도 설정에 따라 결과가 달라지면 두 결과와 바뀐 설정을 함께 남긴다. 두 장면의 연결 근거가 없으면 판정의 범위를 장면 안으로 한정하고, 의도적으로 만든 변형 사례는 규칙 시험 자료로 구분한다. 이런 기록은 검토 순서를 정하고 수정 후 같은 입력을 다시 비교하는 데 사용할 수 있다. 실제 학습 자료의 채택ㆍ삭제나 학습 성능 개선을 측정한 절차는 아니다.

결과의 활용 범위도 층별로 다르다.144개 작성 사례는 조건 연결과 상태 기록이 의도대로 실행되는지, 자연27장면은 원문 해석과 객체ㆍ시점 근거가 어느 곳에서 부족한지 보여 준다.2/7/18 궤적 요약은 문장 모순 판정과 기록 행동 정렬이 별개라는 점을 설명한다.134건 반응 분석과 설정 의존57건은 시간 기준ㆍ계산 증거ㆍ연결 근거를 보존해야 하는 이유를 보여 준다. 따라서 자료를 검토할 때 한 종류의 통과 결과로 다른 종류의 근거 부족을 대신 채우지 않는다.
\end{revision}
